% !TeX program = pdflatex
\documentclass[11pt,a4paper]{amsart}

\newcommand{\notelanguage}{finnish}
% Shared preamble for course notes and exercise sets.

\usepackage[T1]{fontenc}
\usepackage{lmodern}
\providecommand{\notelanguage}{english}
\usepackage[finnish,english]{babel}
\AtBeginDocument{\expandafter\selectlanguage\expandafter{\notelanguage}}

\newcommand{\notesfinnishlanguage}{finnish}
\ifx\notelanguage\notesfinnishlanguage
\newcommand{\notestheoremname}{Lause}
\newcommand{\noteslemmaname}{Lemma}
\newcommand{\notespropositionname}{Propositio}
\newcommand{\notescorollaryname}{Seurauslause}
\newcommand{\notesdefinitionname}{Määritelmä}
\newcommand{\notesexamplename}{Esimerkki}
\newcommand{\notesproblemname}{Tehtävä}
\newcommand{\notesexercisename}{Tehtävä}
\newcommand{\notesremarkname}{Huomautus}
\newcommand{\notessolutionname}{Ratkaisu}
\else
\newcommand{\notestheoremname}{Theorem}
\newcommand{\noteslemmaname}{Lemma}
\newcommand{\notespropositionname}{Proposition}
\newcommand{\notescorollaryname}{Corollary}
\newcommand{\notesdefinitionname}{Definition}
\newcommand{\notesexamplename}{Example}
\newcommand{\notesproblemname}{Problem}
\newcommand{\notesexercisename}{Exercise}
\newcommand{\notesremarkname}{Remark}
\newcommand{\notessolutionname}{Solution}
\fi

\usepackage[a4paper,margin=30mm]{geometry}
\usepackage{microtype}
\usepackage{graphicx}
\usepackage{enumitem}

\newcommand{\exerciseheading}[2]{%
  \par\noindent
  {\large\bfseries #1\par}
  \vspace{0.1em}
  \noindent{\large #2\par}
  \vspace{1.5em}
}

\usepackage{amsmath}
\usepackage{amssymb}
\usepackage{amsthm}
\usepackage{mathtools}
\usepackage{aliascnt}

\usepackage[hidelinks,pdfusetitle]{hyperref}

\numberwithin{equation}{section}
\setlist{itemsep=0.25em,topsep=0.5em}

\theoremstyle{plain}
\newtheorem{theorem}{\notestheoremname}[section]
\newaliascnt{lemma}{theorem}
\newtheorem{lemma}[lemma]{\noteslemmaname}
\aliascntresetthe{lemma}
\newaliascnt{proposition}{theorem}
\newtheorem{proposition}[proposition]{\notespropositionname}
\aliascntresetthe{proposition}
\newaliascnt{corollary}{theorem}
\newtheorem{corollary}[corollary]{\notescorollaryname}
\aliascntresetthe{corollary}

\theoremstyle{definition}
\newaliascnt{definition}{theorem}
\newtheorem{definition}[definition]{\notesdefinitionname}
\aliascntresetthe{definition}
\newaliascnt{example}{theorem}
\newtheorem{example}[example]{\notesexamplename}
\aliascntresetthe{example}
\newaliascnt{problem}{theorem}
\newtheorem{problem}[problem]{\notesproblemname}
\aliascntresetthe{problem}
\newtheorem{exercise}{\notesexercisename}

\theoremstyle{remark}
\newaliascnt{remark}{theorem}
\newtheorem{remark}[remark]{\notesremarkname}
\aliascntresetthe{remark}

\newenvironment{solution}
{
\begin{proof}[\notessolutionname]}
  {
\end{proof}}

\usepackage[nameinlink,noabbrev]{cleveref}

\crefname{theorem}{theorem}{theorems}
\crefname{lemma}{lemma}{lemmas}
\crefname{proposition}{proposition}{propositions}
\crefname{corollary}{corollary}{corollaries}
\crefname{definition}{definition}{definitions}
\crefname{example}{example}{examples}
\crefname{problem}{problem}{problems}
\crefname{exercise}{exercise}{exercises}
\crefname{remark}{remark}{remarks}

\ifx\notelanguage\notesfinnishlanguage
\crefname{theorem}{lause}{lauseet}
\crefname{lemma}{lemma}{lemmat}
\crefname{proposition}{propositio}{propositiot}
\crefname{corollary}{seurauslause}{seurauslauseet}
\crefname{definition}{määritelmä}{määritelmät}
\crefname{example}{esimerkki}{esimerkit}
\crefname{problem}{tehtävä}{tehtävät}
\crefname{exercise}{tehtävä}{tehtävät}
\crefname{remark}{huomautus}{huomautukset}
\fi

\newcommand{\NN}{\mathbb{N}}
\newcommand{\NNone}{\mathbb{N}_1}
\newcommand{\ZZ}{\mathbb{Z}}
\newcommand{\QQ}{\mathbb{Q}}
\newcommand{\RR}{\mathbb{R}}
\newcommand{\CC}{\mathbb{C}}
\newcommand{\PP}{\mathbb{P}}

\DeclareMathOperator{\im}{im}
\DeclareMathOperator{\rank}{rank}
\DeclarePairedDelimiter{\abs}{\lvert}{\rvert}
\DeclarePairedDelimiter{\norm}{\lVert}{\rVert}
\newcommand{\set}[1]{\left\lbrace#1\right\rbrace}
\newcommand{\maxset}[1]{\max\set{#1}}
\newcommand{\minset}[1]{\min\set{#1}}
\newcommand{\powerset}[1]{\mathcal{P}\!\left(#1\right)}


\title{Johdatus yliopistomatematiikkaan\\Harjoituksen 5 ratkaisut}
\date{}
\author{}

\begin{document}

\exerciseheading{Johdatus yliopistomatematiikkaan}{Harjoituksen 5 ratkaisut}

\begin{exercise}
  Osoita, että äärettömien bittijonojen joukko on ylinumeroituva.
\end{exercise}

\begin{exercise}
  Eräässä 150 henkilön ryhmässä 45 harrastaa uintia, 40 pyöräilyä ja 50 juoksua.
  Lisäksi 32 henkilöä harrastaa juoksua mutta ei pyöräilyä, 27 henkilöä
  harrastaa juoksua ja uintia ja 10 henkilöä harrastaa kaikkia kolmea lajia.
\begin{enumerate}[label=\alph*), nosep]
  \item Kuinka moni henkilöistä harrastaa juoksua mutta ei uintia
    eikä pyöräilyä?
  \item Jos 21 henkilöä harrastaa sekä pyöräilyä että uintia, niin kuinka
    moni ei harrasta mitään kyseisistä kolmesta lajista?
\end{enumerate}
Käytä apuna Vennin kaavioita.
\end{exercise}

\begin{exercise}
Yleistä summaperiaate \(n\) joukolle ja todista se induktiolla.
\end{exercise}

\begin{exercise}
Kuinka monessa \(n\) bitin jonossa on
\begin{enumerate}[label=\alph*), nosep]
\item kolme ensimmäistä bittiä ykkösiä ja muut nollia,
\item kolme peräkkäistä ykköstä ja muut nollia?
\end{enumerate}
\end{exercise}

\begin{exercise}
Tarkastellaan kompleksilukuja \(z = (2, -3)\) ja \(w = (4, 1)\).
\begin{enumerate}[label=\alph*), nosep]
\item Esitä \(z\) ja \(w\) imaginaariyksikön avulla muodossa \(x + iy\).
\item Mikä on kompleksiluvun \(z\) liittoluku?
\item Mikä on kompleksiluvun \(w\) vastaluku?
\item Laske kompleksilukujen \(z\) ja \(w\) summa ja erotus.
\item Mikä on kompleksiluvun \(z\) itseisarvo eli moduli?
\end{enumerate}
Esitä määrittämäsi kompleksiluvut piirtämällä niitä vastaavat
kompleksitason pisteet koordinaatistoon.
\end{exercise}

\begin{exercise}
Osoita, että kompleksilukujen kertolasku on \emph{vaihdannainen} eli
\(z \cdot w = w \cdot z\) kaikilla \(z, w \in \CC\).
Mikä on kompleksilukujen kertolaskun \emph{neutraalialkio} eli sellainen
kompleksiluku \(a\), jolla \(a \cdot z = z \cdot a = z\) kaikilla
\(z \in \CC\)?
\end{exercise}

\end{document}
